\section{Мета роботи та завдання}

\textbf{Мета роботи:}
\begin{enumerate}
    \item Ознайомитись з основами написання класів;
    \item Освоїти інкапсуляцію класів;
    \item Розібратися із використанням конструктора;
    \item Навчитися використовувати модифікатори доступу.
\end{enumerate}

\skipline
\textbf{Завдання на лабораторну роботу:}
\begin{enumerate}
    \item \textbf{Створіть клас «Товар» (\texttt{Item}).}
    \begin{description}
        \item[А.] 
        Клас «Товар» має містити такі поля: найменування (рядок), ціна (\texttt{float}).
        \item[Б.] 
        Клас «Товар» повинен містити конструктор, який приймає два параметри: найменування 
        та початкову ціну товару.
        \item[В.] 
        Клас «Товар» повинен мати такі публічні методи: підвищення ціни на певний відсоток 
        (значення відсотка типу \texttt{float} передається як аргумент методу); зниження ціни 
        на певний відсоток (значення відсотка типу \texttt{float} передається як аргумент методу).
        \item[Г.] У класі має бути реалізована перевірка ціни на негативне значення. 
        Якщо в конструкторі передається негативне значення ціни, чи в результаті зниження 
        ціни на $>100\%$ ціна стає негативною, вона має бути примусово встановлена в 0.
    \end{description}

    \item \textbf{Створіть клас «Кошик» (\texttt{Cart}).}
    \begin{description}
        \item[А.] «Кошик» має реалізовувати структуру даних стек, у якому містяться об'єкти 
        класу «Товар». Вважатимемо, що ми додаємо завжди по 1 одиниці товару. Стек «всередині» 
        реалізується традиційним масивом об'єктів \texttt{Item}. Клас повинен містити перевірки, 
        пов'язані з роботою стека — переповнення стека, спроба витягти елемент із порожнього стека 
        і т.д.
        \item[Б.] «Кошик» повинен містити конструктор з 1 параметром — максимальною кількістю 
        елементів у стеку. У цьому конструкторі повинна відбуватися ініціалізація масиву, який 
        реалізує стек.
        \item[В.] «Кошик» повинен містити такі публічні методи: додавання товару, видалення 
        товару, підрахунок суми цін товарів у кошику (пройтися за елементами масиву та скласти 
        всі значення цін), підвищення та зниження цін усіх товарів у стеку (два окремі методи, 
        значення відсотка передається як параметр методу, необхідно пройтись по всіх елементах 
        масиву та передати відповідне повідомлення об'єктам).
    \end{description}

    \item \textbf{У методі \texttt{Main} необхідно створити об'єкт класу «Кошик» з деякою максимальною 
    кількістю елементів у стеку.}
    \begin{description}
        \item[А.] Заповнити кошик об'єктами класу \texttt{Item} (5--6 об'єктів буде достатньо);
        \item[Б.] Вивести суму цін товарів усередині кошика;
        \item[В.] Підняти ціни в кошику на $15\%$, вивести змінену суму цін на консоль;
        \item[Г.] Знизити ціни в кошику на $30\%$, вивести змінену суму цін на консоль.
    \end{description}

    \item \textbf{Перепишіть клас \texttt{Cart}, використовуючи замість стека чергу 
    (або колекцію, якщо ви знаєте, що це таке) без зміни інтерфейсу класу. Перевірте роботу коду.}

    \item \textbf{Додаткове індивідуальне завдання з файлу лаб2-2.docx. Варіант №7}\\
    Облік товарів на складі — Створіть просту консольну програму для обліку товарів. 
    Програма повинна працювати з масивом назв товарів та відповідним масивом їхньої кількості. 
    Необхідно вивести всі товари, визначити товар із найбільшою кількістю одиниць, підрахувати 
    загальну кількість товарів на складі та знайти товари, кількість яких менша за задане 
    користувачем мінімальне значення. Для перебору масивів використайте цикли, а для перевірки 
    кількості — умовні конструкції. 
\end{enumerate}
